\begin{helper}[Terminal PE whose support appears earlier]
Let \(\mathcal E\) be a positive-rung maximum ETT on a finite loopless
multigraph, with displayed tree \(T\), certified series \(\mathcal T\), and
rung number \(n\).  Assume lower-rung assertion (v) for every positive-rung
maximum ETT on the same input with fewer than \(n\) rungs: under a proper
exactly-\(e_0\)-uncolored stable coloring of its terminal ladder tree, a
nonminimum boundary vertex of a defective color precedes every later PE
support or SE extension vertex, from its first-appearance slot onward.

Let \(a=\operatorname{Fin.last}(n)\), suppose the data at \(a\) certify a PE
with chosen support \(v\), let \(h:v\in V(T)\), and suppose
\[
 \noderef{ETTFirstAppearanceIndex}(\mathcal E,v,h)=b+1,
 \qquad b+1<n.
\]
If \(\sigma\) is proper, leaves exactly \(e_0\) uncolored, and is stable on
the terminal ladder tree relative to \((D_a,\varphi_a)\), then \(\sigma\)
is \(\noderef{StableColoringBeforeVertex}\) before \(v\) in \(T\), relative
to \((D_b,\varphi_b)\).  If \(b>0\), it is also
\(\noderef{StableColoring}\) on the ladder tree at slot \(b\) relative to
that pair.
\end{helper}
\begin{proof}
Put \(T^*=T(v)-v\).  The terminal PE certificate supplied by
\noderef{TashkinovSeries}, its ordered-prefix clause in
\noderef{TashkinovSeriesStep}, and
\noderef{TreeSequenceInitialSegment} put every vertex of \(T^*\) in the
terminal ladder tree \(T_n\).  Repeatedly applying
\noderef{HkuMissingStoredMonotone} from slot \(b\) through slot \(n\) gives
\[
 \overline\varphi_b(T_{b+1})\cup D_b
 \subseteq \overline\varphi_n(T_n)\cup D_n.                    \tag{1}
\]
Using a vertex witness from \noderef{TreeMissingColors} and an endpoint
witness from \noderef{TreeIncidentEdge}, the two clauses of the given
stability relation restrict to \(T^*\) with control \(D_b\).  Hence
\(\sigma\) is stable from \(\varphi_n\) there.  By \noderef{sc2}, it remains
to prove that \(\varphi_n\) is stable from \(\varphi_b\) on \(T^*\) with
control \(D_b\).

Eliminate the terminal PE certificate.  Its
\noderef{MaximumDefectiveChoice} supplies a proper exactly-\(e_0\)-uncolored
selected coloring \(\pi_{n-1}\), stable from \(\varphi_{n-1}\) on \(T_n\),
a defective supporting vertex \(v\), a selected boundary-edge identity
\(f_n\) incident with \(v\), an outside endpoint, and
\(\pi_{n-1}(f_n)=\delta_n\).  Since \(b<n\), the complete stage where \(v\)
first appears is contained in \(T_{n-1}\), so \(v\in V(T_{n-1})\).

We verify the lower-(v) hypotheses using exact edge identities.  By
\noderef{DefectiveEdge}, \noderef{BoundaryColorEdges}, and
\noderef{DefectiveColor}, there is a second identity \(g\ne f_n\) in the
\(\delta_n\)-colored boundary of \(T_n\).  Expand
\noderef{BoundaryEdges} to write its endpoints as an inside vertex \(w\)
and an outside vertex.  We have \(w\ne v\): otherwise the distinct edges
\(f_n,g\) are adjacent at \(v\) and carry the same color, contrary to
\noderef{IsProperEdgeColoring}.  The predicates
\noderef{DefectiveEdge} and \noderef{DefectiveVertex} make \(w\) another
defective vertex for the same admissible coloring.  Maximality in
\noderef{MaximumDefectiveChoice} gives \(w=v\) or \(w\prec v\), so
\(w\prec v\).

The first-appearance equality and \(b<n\), expanded through
\noderef{ETTFirstAppearanceIndex}, put \(v\) in \(T_{b+1}\subseteq
T_{n-1}\).  Because \(w\prec v\), the ordered-prefix equalities put \(w\)
there as well.  Both outside endpoints remain outside the smaller tree.
Thus \(f_n\) and \(g\) are two distinct \(\delta_n\)-colored boundary
identities of \(T_{n-1}\).  Therefore \(\delta_n\) is defective there, and
\noderef{BoundaryColorVertex} puts both \(v\) and \(w\) in its boundary
incidence set, with \(w\prec v\).  Apply
\noderef{CertifiedETTRestriction} at rung \(n-1\); its inherited maximum
property and the assumed lower assertion
(v) consequently imply
\[
 v\preceq v_t
 \quad\hbox{for every PE supporting vertex }v_t
 \hbox{ with }b+1\le t\le n-1.                                \tag{2}
\]
All trees, colorings, computed stored sets, kinds, vertices, and relevant
certificates used here are the retained data of that restriction.

Every \(x\in V(T^*)\) satisfies \(x\prec v\).  Equation (2) says that no PE
in the interval from \(b+1\) through \(n\) uses \(x\) as its support.
Applying \noderef{HkuUnsupportedVertexColors} interval by interval yields
\[
 \overline\varphi_b(x)=\overline\varphi_n(x)
 \qquad(x\in V(T^*)).                                         \tag{3}
\]
For the edge clause, take an edge \(f\) incident with \(T^*\), with
\(\varphi_b(f)=c\) and
\(c\in\overline\varphi_b(T^*)\cup D_b\).
\noderef{SameColorSeriesPreserved} says that either
\(\varphi_n(f)=\varphi_b(f)\), or \(f\) is the exact connecting edge of a
PE in this interval.  In the exceptional case the unique tree endpoint is
that PE support.  Since \(f\) is incident with \(T^*\), this support would
strictly precede \(v\), contradicting (2); it cannot equal \(v\) because
\(T^*\) excludes \(v\).  Hence
\[
 \varphi_n(f)=\varphi_b(f).                                    \tag{4}
\]
Equations (3) and (4) are exactly the vertex and edge clauses of stability
from \(\varphi_b\) to \(\varphi_n\) on \(T^*\).  Compose it with the
restricted given stability using \noderef{sc2}; this proves the
before-support conclusion.

Finally assume \(b>0\).  Minimality in
\noderef{ETTFirstAppearanceIndex} says that \(v\) is absent from the ladder
tree at slot \(b\) but present at the next complete stage.  The certified
step at that index and \noderef{TreeSequenceInitialSegment} show, entry by
entry, that every vertex of the predecessor ladder tree occurs before
\(v\) in the final display; injectivity in \noderef{TreeSequence} rules out
an alternative display position.  Thus its vertex set is contained in
\(V(T^*)\).  The witness definition \noderef{TreeIncidentEdge} transports
edge incidence to \(T^*\), and a vertex witness transports missing colors.
Restricting the already proved two stability clauses to this subset gives
stability on the predecessor ladder tree.  No such claim is made when
\(b=0\), exactly as required.
\end{proof}
